% \clearpage
\section{SparseAD的变体范式}
\label{appendix-variant}
\subsection{动机与方法}

\subsubsection{动机}

如 \cref{sec:sparse_perception} 所述，检测查询和跟踪查询之间存在明显的监督不平衡，导致检测性能不可忽视的下降，同时也影响跟踪性能。为缓解该问题，重新审视了稀疏感知的整个框架，认识到场景级检测记忆查询（同样源自历史帧）在一定程度上可能能够充当跟踪查询的角色。此外，在训练阶段，场景级检测记忆查询不通过ID关系与真实标签强制绑定，因此不存在监督不平衡的问题。换句话说，可能仅使用场景级检测记忆查询即可实现端到端跟踪，而不会导致检测性能下降。

\subsubsection{方法}

在训练阶段，丢弃来自实例级记忆的跟踪查询，从而可以省略 \cref{sec:train_infer_detail} 中提到的\textbf{第二阶段}。同时，将可学习检测查询与场景级检测记忆查询的比例从原来的644:256调整为200:700，这有助于保证跟踪的一致性。为替代跟踪查询并为下游任务提供稳定的障碍物历史信息，将匹配同一真实标签的检测查询拼接形成一个新的实例。也就是说，一个实例的历史信息可能来自不同的检测查询。在推理阶段，场景级检测记忆查询对应的边界框将继承前一帧的ID，而可学习检测查询对应的边界框将被分配新的ID。与训练阶段相同，将具有相同ID的查询拼接形成新的实例供下游任务使用。该范式的其他细节与 \cref{sec:method} 和 \cref{appendix-imple} 中的描述保持一致。

\subsection{验证实验}
在nuScenes val集上进行实验。多个驾驶任务的性能如 \cref{tab:hard-asso-det-tracking} 所示。正如预期，无硬关联的\model 在感知性能上取得了显著提升。与SparseAD-Base相比，NDS提升了 $+1.2\%$，mAP提升了 $+2.4\%$，AMOTA提升了 $+1.7\%$。使用ViT-Large作为图像骨干网络时，NDS提升了 $+2.3\%$，mAP提升了 $+3.7\%$，AMOTA提升了 $+2.3\%$。然而，需要指出的是，由于训练过程中缺乏硬关联，IDS也出现了明显增加。同时比较了两种方法的预测和规划性能，结果如 \cref{tab:hard-asso-motion-planner} 所示。尽管无硬关联的\model 在某些指标上感知表现更好，但两种方法的预测和规划能力相当。本文认为，虽然无硬关联的\model 召回了更多困难样本，但过高的IDS意味着轨迹一致性不够好，最终导致预测和规划性能相似。

\begin{table*}[t]
		\caption{硬关联监督对SparseAD稀疏感知性能影响的对比实验。}
		\begin{center}
			\tiny
			\vspace{-10pt}
			\resizebox{1\textwidth}{!}
			{
				\begin{tabular}{c|c|cc|ccccc|ccccc|c}
					\toprule
					\multirow{2}{*}{骨干网络}& \multirow{2}{*}{硬关联}& \multicolumn{7}{c|}{检测} & \multicolumn{5}{c}{跟踪}& 建图\\
					&&NDS$\uparrow$&	mAP$\uparrow$&	ATE$\downarrow$&	ASE$\downarrow$&	AOE$\downarrow$	&AVE$\downarrow$&	AAE$\downarrow$& AMOTA$\uparrow$&	AMOTP$\downarrow$&	Recall$\uparrow$ & IDs$\downarrow$ &Frag$\downarrow$& mAP$\uparrow$\\
					\midrule
					\multirow{2}{*}{V2-99\cite{lee2020centermask}}& \Checkmark& 0.579&	0.475&	0.556&	0.264&	0.295&	0.288&	0.186&0.530&	1.087&	0.608 &297&	454& 34.2\\
					&\XSolidBrush &0.591&	0.499&	0.554&	0.260&	0.323&	0.256&	0.198& 0.547&	1.115&	0.629 &639&	658	& 33.5\\
					\midrule
					\multicolumn{2}{|c|}{\textit{差值}}& \textcolor{red}{+0.012}& \textcolor{red}{+0.024} &\textcolor{red}{-0.002} &\textcolor{red}{-0.004}& \textcolor{blue}{+0.028} &\textcolor{red}{-0.032}& \textcolor{blue}{+0.008}& \textcolor{red}{+0.017} &\textcolor{blue}{+0.028}& \textcolor{red}{+0.021}& \textcolor{blue}{+442}& \textcolor{blue}{+204}&\textcolor{blue}{-0.7}\\
					\midrule
					\multirow{2}{*}{ViT-L\cite{dosovitskiy2020image}}& \Checkmark& 0.625&	0.536&	0.531&	0.252&	0.219&	0.230&	0.199& 0.606&	0.933&	0.706&251&	335	& 34.6\\
					&\XSolidBrush&0.648&	0.573&	0.500&	0.251&	0.223&	0.210&	0.202&0.629&	0.942&	0.733 &958&	507 & 34.3\\
					\midrule
					\multicolumn{2}{|c|}{\textit{差值}} & \textcolor{red}{+0.023}& \textcolor{red}{+0.037}& \textcolor{red}{-0.031}& \textcolor{red}{-0.001}& \textcolor{blue}{+0.004}& \textcolor{red}{-0.020}& \textcolor{blue}{+0.003}& \textcolor{red}{+0.023}& \textcolor{blue}{+0.009}& \textcolor{red}{+0.027}& \textcolor{blue}{+707}& \textcolor{blue}{+172} & \textcolor{blue}{-0.3}\\
					\bottomrule
				\end{tabular}% }
			}
		\end{center}
		\label{tab:hard-asso-det-tracking}
	\end{table*}

	\begin{table*}[t]
		\caption{硬关联监督对SparseAD运动规划器性能影响的对比实验。}
		\begin{center}
			\tiny
			\vspace{-10pt}
			{
				\begin{tabular}{c|c|ccc|cccc|cccc}
					\toprule
					\multirow{2}{*}{骨干网络}& \multirow{2}{*}{硬关联}& \multicolumn{3}{c|}{运动预测} & \multicolumn{4}{c|}{规划L2(m)$\downarrow$} & \multicolumn{4}{c}{规划碰撞率$\downarrow$}\\
					&&minADE$\downarrow$&	minFDE$\downarrow$&	MR$\downarrow$&1s &2s&3s&平均&1s &2s&3s&平均\\
					\midrule
					\multirow{2}{*}{V2-99\cite{lee2020centermask}}& \Checkmark& 0.830&	1.584&	0.187& 0.15&0.31&0.57&0.35 &	0.00&0.06&0.21&0.09\\
					&\XSolidBrush & 0.823&	1.551&	0.191&	0.15&0.31&0.56&0.34&	0.00&0.04&0.23& 0.09\\
					\midrule
					\multicolumn{2}{|c|}{\textit{差值}}& \textcolor{blue}{+0.003} &\textcolor{red}{-0.033}& \textcolor{blue}{+0.004}& 0.00& 0.00& \textcolor{red}{-0.01}& \textcolor{red}{-0.01}& 0.00 &\textcolor{red}{-0.02}& \textcolor{blue}{+0.02}& 0.00\\
					\midrule
					\multirow{2}{*}{ViT-L\cite{dosovitskiy2020image}}& \Checkmark&0.741&	1.458&	0.169& {0.15}&{0.31}&{0.56}&{0.34}&	{0.00}&{0.04}&{0.15}&{0.06}\\
					&\XSolidBrush& 0.740&	1.453&	0.173&	0.15&0.31&0.56&0.34&	0.00&0.04&0.14& 0.06\\
					\midrule
					\multicolumn{2}{|c|}{\textit{差值}} & \textcolor{red}{-0.001}& \textcolor{red}{-0.005}& \textcolor{blue}{+0.004}& 0.00& 0.00 &0.00& 0.00& 0.00& 0.00& \textcolor{red}{-0.01} &0.00\\
					\bottomrule
				\end{tabular}% }
			}
		\end{center}
		\label{tab:hard-asso-motion-planner}
	\end{table*}
